\section{Where the carbon and money go}
\label{sec:ledgers}
The headline result is easier to audit as two ledgers. A downward arrow means the candidate causes less than the conventional baseline; an upward arrow means the candidate introduces a new lifecycle burden or cost. These rows are generated from the same Rust functions as the final result, so the tables cannot be reconciled by manually changing manuscript numbers.

\subsection{Annual hydrogen production}
INL reports the common hydrogen service as 130 MMSCFD and approximately 29,000 lb/h~\cite{inltev953,inltev961}. Using Assumption A-02, $A=0.85$, the mass-rate annualisation is
\begin{equation}\label{eq:annual-h2}
M_{H_2}=29{,}000\frac{\mathrm{lb}}{\mathrm h}
\left(0.45359237\frac{\mathrm{kg}}{\mathrm{lb}}\right)
(8760\,\mathrm{h/y})(0.85)
\left(\frac{1\,\mathrm t}{1000\,\mathrm{kg}}\right)
=97{,}946~\mathrm{t/y}.
\end{equation}
This is a project annualisation of the source instantaneous rate, not a separate literature-reported annual plant output.

\subsection{Where did the CO$_2$ go?}
The source process reports 3,205 short ton/day baseline emissions~\cite{inltev953}, and 142 short ton/day residual emissions plus 1,927 short ton/day captured CO$_2$ for Case 6~\cite{inltev961}. These source values are converted rather than reverse-engineered from an assumed pure-methane feed, because the INL process model already distinguishes feed, fuel and process streams. The annualisation uses the Nishihara design-study availability $A=0.85$ (Assumption A-02)~\cite{nishihara2007potential}. Using the standard conversion 1 short ton = 0.90718474 t:
\begin{equation}\label{eq:base-direct}
E_{\rm base,direct}=3205\frac{\mathrm{short\ ton}}{\mathrm d}
(0.90718474)\frac{\mathrm t}{\mathrm{short\ ton}}(365)(0.85)
=902{,}060~\mathrm{t/y}.
\end{equation}
\begin{equation}\label{eq:cand-direct}
E_{\rm cand,direct}=142(0.90718474)(365)(0.85)
=39{,}966~\mathrm{t/y}.
\end{equation}
so
\begin{equation}\label{eq:direct-avoided}
A_{\rm direct}=902{,}060-39{,}966=862{,}094~\mathrm{tCO_2/y}.
\end{equation}
The same conversion gives
\begin{equation}\label{eq:captured-annual}
M_{\rm captured}=1927(0.90718474)(365)(0.85)
=542{,}362~\mathrm{tCO_2/y}.
\end{equation}

Natural-gas use falls from 52.5 MMSCFD~\cite{inltev953} to 34.0 MMSCFD~\cite{inltev961}. Assumption A-03 uses 1044 Btu/scf as the project screening HHV for volumetric-to-energy conversion; it is not asserted as an INL-measured gas property. Assumption A-04 uses the IEA global-average 11.5 kgCO$_2$e/GJ gas-supply factor~\cite{iea2026lng}. Thus the upstream lifecycle terms are approximately 206,322 and 133,618 tCO$_2$e/y, respectively, saving 72,704 tCO$_2$e/y.  The new candidate burdens use Assumption A-05, the UNECE 5.5 kgCO$_2$e/MWh$_e$ nuclear-electricity lifecycle anchor~\cite{unece2022lca}. For direct heat, Assumption A-06 multiplies that electricity factor by the declared 0.504 MWh$_e$-equivalent/MWh$_{th}$ allocation proxy; this is not a published direct-heat LCA. The auxiliary term uses the source 17.3 MWe candidate demand~\cite{inltev961} minus Assumption A-01 = 6.3 MWe, a retained historical project screening anchor rather than an INL common-output value. Assumption A-07 = 0.025 tCO$_2$e/tCO$_2$ is likewise a generic project T\&S lifecycle screen, not a Singapore-route measurement. Therefore
\begin{equation}\label{eq:nuclear-lca}
E_{\rm nuclear}=176.8(1000)(7446)(5.5\times0.504)/10^6
=3{,}649~\mathrm{tCO_2e/y},
\end{equation}
\begin{equation}\label{eq:aux-lca}
E_{\rm aux}=(17.3-6.3)(1000)(7446)(5.5)/10^6
=450~\mathrm{tCO_2e/y},
\end{equation}
and
\begin{equation}\label{eq:ts-lca}
E_{\rm T\&S}=0.025(542{,}362)=13{,}559~\mathrm{tCO_2e/y}.
\end{equation}
Direct CO$_2$ and lifecycle CO$_2$e are not interchangeable. The conventional and candidate plant stacks first determine the direct saving. Lifecycle accounting then adds the change in upstream natural-gas emissions and subtracts new nuclear, auxiliary-electricity and CO$_2$ transport/storage burdens:
\begin{equation}\label{eq:lifecycle-avoided}
A_{\rm lifecycle}=A_{\rm direct}+A_{\rm upstream}-E_{\rm nuclear}-E_{\rm aux}-E_{\rm T\&S}.
\end{equation}

\begin{table}[htbp]\centering\scriptsize
\caption{Canonical lifecycle-emissions ledger generated by Rust. SAVED means less than the conventional baseline; ADDED means a new candidate burden.}
\label{tab:co2-ledger}
\pgfplotstabletypeset[col sep=comma,
columns={item,direction,change_tco2e_y,meaning},
columns/item/.style={string type,column name=Item,column type={p{0.31\linewidth}}},
columns/direction/.style={string type,column name=Direction,column type={p{0.11\linewidth}}},
columns/change_tco2e_y/.style={fixed,precision=0,column name={tCO$_2$e/y},column type=r},
columns/meaning/.style={string type,column name=Meaning,column type={p{0.31\linewidth}}},
every head row/.style={before row=\toprule,after row=\midrule},
every last row/.style={after row=\bottomrule}]
{../results/final_design/generated/lifecycle_ledger.csv}
\end{table}

\begin{figure}[htbp]\centering
\input{../results/final_design/generated/co2_bridge_figure.tex}
\caption{Rust-generated lifecycle contribution bridge.  Positive bars are savings; negative bars are candidate lifecycle burdens.  The net bar reproduces the canonical lifecycle result.}
\label{fig:co2-bridge}
\end{figure}

The arithmetic reconciliation is therefore
\[
862{,}093.8+72{,}703.8-3{,}649.2-450.5-13{,}559.0
=\mathbf{917{,}138.9~tCO_2e/y}.
\]
Lifecycle avoidance exceeds direct avoidance because the approximately 72.7 ktCO$_2$e/y upstream natural-gas saving is larger than the approximately 17.7 ktCO$_2$e/y of new lifecycle burdens. Captured CO$_2$ is not added again to the numerator, avoiding double counting.  The one-tonne-scale difference obtained by adding rounded table entries is rounding only.

\subsection{Where did the money go?}
The natural-gas saving is directly reproducible from the same cited 52.5 and 34.0 MMSCFD source rates. Assumption A-08 sets the gas price to S\$15/GJ as a project screening value, not a contracted Singapore tariff:
\begin{equation}\label{eq:ng-costs}
C_{\rm NG,base}\approx\mathrm{S\$}269.115~\mathrm{million/y},\qquad
C_{\rm NG,cand}\approx\mathrm{S\$}174.284~\mathrm{million/y}.
\end{equation}
hence
\begin{equation}\label{eq:ng-saving}
S_{\rm NG}=269.115-174.284
=\mathbf{S\$94.831~million/y}.
\end{equation}

The selected source economic burden is deliberately charged in full. Nishihara reports 0.52 JPY/MJ heat and 4.9 JPY/kWh electricity at the source design basis~\cite{nishihara2007potential}. The repository normalizes 2007 JPY using the recorded 2007-to-2025 Japan deflator ratio 112.27/99.59=1.12732 and dated FX 0.008117 SGD/JPY, giving approximately S\$4.758/GJ and S\$44.837/MWh. These are broad screening normalizations, not nuclear-specific construction indices. Therefore
\begin{equation}\label{eq:source-heat-cost}
C_{\rm source,heat}=370(7446)(3.6)(4.758)\approx\mathrm{S\$}47.193~\mathrm{million/y},
\end{equation}
\begin{equation}\label{eq:source-electric-cost}
C_{\rm source,electric}=88(7446)(44.837)\approx\mathrm{S\$}29.380~\mathrm{million/y},
\end{equation}
and
\begin{equation}\label{eq:source-total-cost}
C_{\rm source}\approx\mathbf{S\$76.572~million/y}.
\end{equation}
The 88 MWe term is used only to recover the full published source economic burden; it is not a project export claim.

The CCS capital starts from IEAGHG Case 1A's 2014 capture-capital basis and is price/currency normalized and linearly scaled by captured tonnes, giving approximately S\$114.04 million~\cite{ieaghg2017smr}. The linear scaling is a project screening assumption. Assumption A-13 retains the IEAGHG 8\% / 25-year economic basis; with $\mathrm{CRF}(i,n)=i(1+i)^n/[(1+i)^n-1]$, annualisation gives
\begin{equation}\label{eq:ccs-annual-cost}
C_{\rm CCS,annual}\approx\mathbf{S\$10.683~million/y}.
\end{equation}
Assumption A-14 is an explicit 10\% project integration/site capital allowance rather than an observed Singapore cost. Assumption A-15 annualises it at the retained 3\% / 40-year design-study convention. It contributes approximately
\begin{equation}\label{eq:integration-cost}
C_{\rm integration}\approx\mathbf{S\$2.857~million/y}.
\end{equation}
while Assumption A-16 sets T\&S to S\$15/tCO$_2$ as a project screening tariff; official Singapore evidence indicates that actual cross-border CCS costs remain unresolved~\cite{mticcs2025cost}. Thus
\begin{equation}\label{eq:ts-cost}
C_{\rm T\&S}=542{,}362(15)
\approx\mathbf{S\$8.135~million/y}.
\end{equation}
Assumption A-17 sets project electricity revenue to exactly \textbf{S\$0/y} because no validated off-design export model exists; this is a conservative model boundary, not a market-price claim.
The economic ledger follows the same logic. The conventional baseline's natural-gas expenditure is replaced by lower candidate gas use plus the full selected source reactor economic burden, CCS annualisation, integration allowance and T\&S. Project electricity revenue is exactly zero.

\begin{table}[htbp]\centering\scriptsize
\caption{Canonical annual-cost ledger generated by Rust for the zero-electricity-credit model boundary.}
\label{tab:cost-ledger}
\pgfplotstabletypeset[col sep=comma,
columns={item,direction,annual_sgd,meaning},
columns/item/.style={string type,column name=Item,column type={p{0.31\linewidth}}},
columns/direction/.style={string type,column name=Direction,column type={p{0.11\linewidth}}},
columns/annual_sgd/.style={fixed,precision=0,column name={S\$/y},column type=r},
columns/meaning/.style={string type,column name=Meaning,column type={p{0.31\linewidth}}},
every head row/.style={before row=\toprule,after row=\midrule},
every last row/.style={after row=\bottomrule}]
{../results/final_design/generated/cost_ledger.csv}
\end{table}

Total money saved is approximately S\$94.831 million/y.  Total new represented costs are approximately S\$98.247 million/y, so
\begin{equation}\label{eq:annual-incremental-cost}
\Delta C_{\rm annual}
\approx98.247-94.831
=\mathbf{S\$3.416~million/y}.
\end{equation}
The final division is therefore
\begin{equation}\label{eq:abatement-cost-final}
C_{\rm abatement}
=\frac{\Delta C_{\rm annual}}{A_{\rm lifecycle}}
\approx\frac{\mathrm{S\$}3.416\ \mathrm{million/y}}{917{,}139\ \mathrm{tCO_2e/y}}
\approx\mathbf{S\$3.725/\mathrm{tCO_2e}}.
\end{equation}

\begin{figure}[htbp]\centering
\input{../results/final_design/generated/cost_bridge_figure.tex}
\caption{Rust-generated annual cost-contribution bridge.  Natural-gas savings are negative cost; reactor, CCS, integration and transport/storage are added costs.  Electricity revenue is zero by model boundary.}
\label{fig:cost-bridge}
\end{figure}

\subsection{What is not fully priced here?}
The low screening abatement cost is not a turnkey Singapore project price. The current boundary does not fully price project-specific FOAK premium, financing structure, Singapore site/construction premium, licensing and security, detailed EPC/contingency, emergency-planning infrastructure, detailed 176.8 MWth IHX qualification, waste/decommissioning, insurance/liability, schedule risk, or negotiated cross-border CCS infrastructure/contracts. No monetary values are invented for these unresolved items. Their omission is why the economic conclusion is threshold compliance within the declared model, not bankability.